% RESOURCE-FOURIER-DIRECTSUM13, 2026-10-19 00:20 KST
\documentclass{article}
\usepackage{amsmath,amssymb,amsthm}
\newtheorem{theorem}{Theorem}
\begin{document}
\section*{LOSSY-FOURIER-DIRECTSUM13: restricted research note}
Status: elementary restricted-model proof pending independent review; P versus NP open.
\begin{theorem}
Let $Q$ source bits be balanced and pairwise independent under a source distribution.
Any fixed public decoder with at most $t$ adaptive bit probes into an arbitrary
possibly nonlinear/many-to-one/randomized $S$-bit encoding, with error $\epsilon_j$
on coordinate $j$, satisfies
\[
\sum_{j=1}^Q (1-2\epsilon_j)^2
\le \sum_{r=0}^{\min(t,S)}\binom Sr.
\]
\end{theorem}
\noindent Proof: each depth-$t$ adaptive decision tree is a polynomial of
degree at most $t$ in the stored bits. All such decoder outputs belong to
one subspace spanned by degree-at-most-$t$ monomials. The centered source
bits form an orthonormal set. Project onto the decoder subspace, apply
Cauchy--Schwarz and Bessel's inequality. Random coins average polynomials.
For the KPI trace-polynomial product with $k\ge2$ coefficients, all $b2^n$
coordinate bits are pairwise independent, giving
$1+S+\binom S2\ge (1-2\epsilon)^2 b2^n$
at average error $\epsilon<1/2$. This concerns a fixed decoder across
databases, not source-specific public algorithms or general SAT machines.
The earlier global basis-cell theorem recovers complete blocks offline but
does not give online two-probe reconstruction. Computational checks are
preserved in the local research bundle if repository publication fails.
\end{document}
% END_OF_VERIFIED_REPORT::RESOURCE-FOURIER-DIRECTSUM13::P_VS_NP_OPEN